\documentclass[10pt]{article}

\usepackage[utf8]{inputenc}

\usepackage[T1]{fontenc}

\usepackage{amsmath}

\usepackage{amsfonts}

\usepackage{amssymb}

\usepackage[version=4]{mhchem}

\usepackage{stmaryrd}



\begin{document}

деления. Формальная структура цепочки не отличается от структуры цепочки в классич. случае. На практике часто используют так наз. импульсное представление квантового Р. в., к-рое является фурье-преобразованием компонент квантового Р. в. по координатам $\vec{q}_{1}, \ldots, \vec{q}_{s}$. А. В. Солдатов. PACПРЕДЕЛЕНИЯ ВЕКТОР к Л а с с и ч е с к и Й - вектор, характеризующий состояние классической системы тождественных частиц. В статистич. механике постулируется, что состояние классич. системы $N$ тождественных частиц полностью характеризуется симметричной относительно перестановок аргументов частиц $x_{1}=\vec{q}_{1} \vec{p}_{1} ; \ldots ; x_{N}=\vec{q}_{N}, \vec{p}_{N}$ функцией распределения



\[

f\left(x_{1}, \ldots, x_{i}, \ldots, x_{j}, \ldots, x_{N}\right)=f\left(x_{1}, \ldots, x_{j}, \ldots, x_{i}, \ldots, x_{N}\right), \text { (1) }

\]



заданной в фазовом пространстве системы и удовлетворяющей условию нормировки:



\[

\int d \vec{q}_{1} \ldots d \vec{q}_{N} d \vec{p}_{1} \ldots d \vec{p}_{N} f\left(\vec{q}_{1}, \vec{p}_{1} ; \ldots ; \vec{q}_{N}, \vec{p}_{N}\right)=1

\]



Зададим приведен̆ную $s$-частичную функцию распределения $f_{s}\left(x_{1}, \ldots, x_{s}\right), s \leqslant N$, симметричную в силу условия (1) относительно перестановки аргументов частиц $x_{1}, \ldots, x_{s}$ :



\[

f_{S}=\frac{N !}{(N-s) !} \int d x_{s+1} \ldots d x_{N} f\left(x_{1}, \ldots, x_{s}, x_{s+1}, \ldots, x_{N}\right),

\]



причем $f_{0} \equiv 1$ в силу нормировки (2), а условие нормировки для $f_{s}\left(x_{1}, \ldots, x_{s}\right)$ имеет вид:



\[

\int d x_{1} \ldots d x_{s} f_{s}\left(x_{1}, \ldots, x_{s}\right)=\frac{N !}{(N-s) !} .

\]



Совокупность приведенных функций распределения (3)



\[

\left(f_{0}, f_{1}, \ldots, f_{N}\right)

\]



можно рассматривать как формальный $(N+1)$-мерный вектор, поэтому ее принято называть ве ктором распре делен и я с исте мы. Компоненты вектора распределения не являются независимыми, так как



$f_{r}\left(x_{1}, \ldots, x_{r}\right)=\frac{(N-s) !}{(N-r) !} \int d x_{r+1} \ldots d x_{s} f_{s}\left(x_{1}, \ldots, x_{s}\right), r<s \leqslant N$.



В термодинамич. пределе $N \rightarrow \infty, V \rightarrow \infty, N / V=$ const Р. в. становится бесконечномерным. При этом



\[

\lim _{N \rightarrow \infty} \frac{N !}{(N-s) !} \Rightarrow N^{s}

\]



поэтому условие нормировки в термодинамич. пределе записывают в форме:



\[

\lim N^{-s} \int d x_{1} \ldots d x_{s} f_{s}\left(x_{1}, \ldots, x_{s}\right)=1,

\]



где предел берется при $N \rightarrow \infty, V \rightarrow \infty, N / V=$ const. Удобство представления (5) заключается в том, что на практике для вычисления макроскопич. величин достаточно знать несколько низших компонент Р. в. Это обстоятельство является основой для разработки приближенных методов определения низших компонент Р. в. Временная эволюция Р. в. определяется цепочкой $N+1$ зацепляющихся уравнений Боголюбова - Борна - Грина - Кирквуда - Ивона (ББГКИ-цепочка). Наряду с определением (3) часто используется также определение для приведенной функции распределения, данное Н. Н. Боголюбовым:



\[

f_{s}^{\mathrm{Б}}=\frac{(N-s) ! V^{s}}{N !} f_{s}

\]



A. В. Солдатов.



РАСПРЕДЕЛЕНИЯ ВЕРОЯТНОСТЕЙ ЧАСТИЧНЫЕ ФУНКЦИи - последовательность функций



\[

F(t)=\left\{F_{s}\left(t, x_{1}, \ldots, x_{s}\right)\right\}_{s \geqslant 1},

\]



которая определяется по решениям Лиувиля уравнения



\[

D(t)=\left\{D_{n}\left(t, x_{1}, \ldots, x_{n}\right)\right\}_{n \geqslant 1}

\]



формулой (см. [1]):



\[

\begin{aligned}

& \quad F_{s}\left(t, x_{1}, \ldots, x_{s}\right)= \\

& =\left(\sum_{n=0}^{\infty} \frac{1}{n !} \int d x_{1} \ldots d x_{n} D_{n}\left(t, x_{1}, \ldots, x_{n}\right)\right)^{-1} \times \\

& \times \sum_{n=0}^{\infty} \frac{1}{n !} \int d x_{s+1} \ldots d x_{s+n} D_{s+n}\left(t, x_{1}, \ldots, x_{s+n}\right) .

\end{aligned}

\]



Для системы фиксированного $N$ числа частиц $D(t)$ - однокомпонентная последовательность:



\[

D(t)=\left\{0, \ldots, 0, D_{N}\left(t, x_{1}, \ldots, x_{N}\right), 0, \ldots\right\} .

\]



Р. в. ч. Ф. удовлетворяют Боголюбова уравнениям.



Равновесные Р. в.ч.Ф. определяются по стационарным решениям уравнения Лиувилля - функциям Бо льцмана:



$D_{n}\left(x_{1}, \ldots, x_{n}\right)=\exp \left\{-\beta \sum_{i=1}^{n}\left(\frac{p_{i}^{2}}{2}-\mu\right)-\beta \sum_{i<j=1}^{N} \Phi\left(q_{i}-q_{j}\right)\right\}$,



где $\mu, \beta$ - параметры равновесного состояния, $\Phi$ - потенциал взаимодействия частиц. Такие функции удовлетворяют уравнениям для равновесных состояний Кирквуда Зальцбурга, Боголюбова, Майера - Монтролла и т. п.



Р. в. ч. ф. $F_{s}\left(t, x_{1}, \ldots, x_{s}\right)$ интерпретируются как плотность вероятности (по лебеговой мере) нахождения $s$ различных частиц в фазовых точках $x_{1}, \ldots, x_{s}$ независимо от фазовых состояний остальных частиц.



Для квантовых систем Р. в. ч. ф. определяются аналогичным образом по решениям уравнений для плотности матрицы.



В терминах Р. в. ч. ф. в статистич. механике описываются состояния систем бесконечного числа частиц.



Лum.: [1] ПетринаД.Я., ГерасименкоВ.И., Малыш е в П. В., Математические основы классической статистической механики, Киев, $1985 . \quad$ В. И. Герасименко.



РАСПРЕДЕЛЕНИЯ ЗАКОН - УдОбНЫЙ опИсатеЛЬНЫЙ термин, могуший означать, в зависимости от контекста, распределение вероятностей (напр.. какой-либо случайной величины) или соответствующую распределения функцию или плотность вероятности. $\quad$ В. И. Битюцков. РАСПРЕДЕЛЕНИЯ ФУНКЦИЯ в клас с иче с коЙ статистической механике - соответствие между макроскопическими наблюдаемыми величинами $A(\vec{x}, t)$, являюшимися полями в физическом пространстве-времени $(\vec{x}, t)$ и соответствуюшими им микроскопическими динамическими переменными $a(p, q, \vec{x}, t)$, которые зависят от координат фазового пространства системы $q=\vec{q}_{1}, \ldots, \vec{q}_{N}$; $p=\vec{p}_{1}, \ldots, \vec{p}_{N} ; N-$ число частиц в системе. Это соответствие задает отображение пространства микросостояний системы на физич. пространство макроскопич. наблюдаемых, поэтому величину $A(\vec{x}, t)$ можно трактовать как функционал от $a(p, q, \vec{x}, t)$



\[

A(\vec{x}, t) \equiv\langle a(p, q, \vec{x}, t)\rangle .

\]



На функционал (1) накладывают естественные ограничения:



\[

\begin{gathered}

\left\langle c_{1} a_{1}(p, q, \vec{x}, t)+c_{2} a_{2}(p, q, \vec{x}, t)\right\rangle=, \\

=c_{1}\left\langle a_{1}(p, q, \vec{x}, t)\right\rangle+c_{2}\left\langle a_{2}(p, q, \vec{x}, t)\right\rangle, \\

\langle c\rangle=c,

\end{gathered}

\]



где $c, c_{1}, c_{2}-$ произвольные постоянные. Функционал, удовлетворяюший условиям (2)-(3), можно представить в виде:



\[

A(\vec{x}, t)=\int d q d p a(p, q, \vec{x}, t) f(q, p) .

\]



Функция $f(q, p)$ определена в фазовом пространстве системы и в силу (3) удовлетворяет условию



\[

\int d q d p f(q, p)=1 .

\]



Функции $f(q, p)$, пригодные для построения функционала (4), называются функциями распределения в фазовом пространстве. Постулируется, что макро-





\end{document}